\documentclass[12pt]{article}
\usepackage[utf8]{inputenc}
\usepackage[T1]{fontenc}
\usepackage[margin=1in]{geometry}
\usepackage{setspace}
\usepackage{xcolor}
\usepackage[colorlinks=true, urlcolor=blue!60!black, linkcolor=blue!60!black]{hyperref}
\setlength{\parindent}{2em}
\setlength{\parskip}{1.25em}
\renewcommand{\baselinestretch}{1.0}
\begin{document}
\begin{center}
{\LARGE \textit{Corfield v. Coryell}, 6 F. Cas. 546 (C.C.E.D. Pa. 1823)\par}
\vspace{0.8em}
{\large Case Brief --- Four-Page Edition\par}
\end{center}
\vspace{1em}

\section{Memory Jogger}

Riding circuit, Justice Bushrod Washington upheld New Jersey's ban on nonresidents harvesting oysters in state waters, famously defining Article IV's Privileges and Immunities Clause as protecting only ``fundamental'' rights belonging to citizens of all free governments---not a cotenancy in a state's common property (\href{https://press-pubs.uchicago.edu/founders/documents/a4_2_1s18.html}{Founders' Constitution}).

\section{Detailed Case Facts}

Corfield owned the vessel \textit{Hiram}, which he leased to Hand, who subleased it to John Keene for \$10 a month. On May 15, 1821, Keene and a crewmate named Courtney dredged for oysters in Maurice River Cove, an indentation of Delaware Bay off southern New Jersey. A New Jersey vessel carrying 40 to 50 armed men, including county magistrates and constables, seized the \textit{Hiram}; defendant Coryell acted as prize master. Two Cumberland County justices of the peace condemned the vessel, which was auctioned for \$10. Coryell relied on New Jersey's 1820 oyster statute, whose Section 6 barred anyone not an actual inhabitant and resident of the state from gathering clams or oysters in New Jersey waters aboard a vessel not wholly owned by a resident, on pain of forfeiture of the vessel.

\section{Procedural History}

Corfield sued Coryell in trespass in the federal Circuit Court for the Eastern District of Pennsylvania. At the April Term 1823 trial, the jury returned a \$560 verdict for Corfield, subject to reserved questions of law. After argument in October 1824, Justice Washington delivered his opinion in April 1825 and ordered judgment for Coryell. No appeal was taken.

\section{Judicial Votes}

\begin{itemize}
\item \textbf{Single opinion:} Associate Justice Bushrod Washington, sitting as Circuit Justice, delivered the sole opinion of the court.
\item \textbf{Concurrences and dissents:} None.
\end{itemize}

\section{Holding}

Judgment for Coryell, setting aside the provisional verdict. (1) New Jersey's statute violates neither the Commerce Clause, the Privileges and Immunities Clause, nor the Admiralty Clause. (2) The Privileges and Immunities Clause protects only rights that are fundamental in their nature and does not give citizens of other states a cotenancy in common property, such as oyster beds, that a state holds for its own citizens. (3) An owner who has leased his vessel for an unexpired term lacks possession and cannot sue in trespass. Because that pleading defect alone defeated the suit, the constitutional rulings are arguably alternative holdings or dicta---and the famous catalogue of fundamental rights was plainly dicta.

\section{Analysis of Opinions}

\textbf{Commerce Clause.} Relying on \textit{Gibbons v. Ogden}, Washington held that Congress's power over navigation and commercial intercourse did not cede to the federal government a state's territory or its proprietary interest in the soil and fisheries beneath its waters. Like inspection laws, the statute regulated removal of oysters from state beds, not commerce in them.

\textbf{Privileges and Immunities.} Washington rejected the claim that nonresidents must share every advantage of local citizens. He confined the Clause to fundamental rights, grouped under protection by government, life and liberty, property, and the pursuit of happiness and safety, and listed examples: passing through or residing in another state for trade or professional pursuits, habeas corpus, suing in state courts, holding property, equal taxation, and the elective franchise as regulated by state law. Fishing in state waters, by contrast, belongs to the state's citizens as ``tenants in common,'' and opening the oyster beds to all comers would risk exhausting a finite resource.

\textbf{Admiralty and form of action.} The states never surrendered authority over their internal fisheries by granting admiralty jurisdiction, and Congress had passed no statute on oyster dredging. After reviewing colonial grants, Washington found it unnecessary to decide the county-boundary question, because Corfield, having leased the \textit{Hiram}, lacked possession to sue in trespass.

\section{Examples: Future Applications}

\begin{enumerate}
\item \textbf{Same-Side: Elk-Hunting Fees.} Montana charges nonresident recreational hunters \$225 and residents \$9. Recreational hunting is not fundamental (\textit{Baldwin v. Fish \& Game Commission}, 1978).
\item \textbf{Same-Side: Citizens-Only Records Access.} Virginia limits its Freedom of Information Act to its own citizens. Statutory records access is not a fundamental right (\textit{McBurney v. Young}, 2013).
\item \textbf{Opposite-Side: Shrimp-Boat Licenses.} South Carolina charges nonresident commercial shrimpers \$2,500 and residents \$25. Pursuing a livelihood is a fundamental privilege, and the fee violates Article IV (\textit{Toomer v. Witsell}, 1948).
\item \textbf{Opposite-Side: Bar Residency.} New Hampshire refuses to admit a qualified Vermont resident to its bar. Professional pursuits are protected (\textit{Supreme Court of New Hampshire v. Piper}, 1985).
\item \textbf{Fence-Sitter: Local Hiring on City Projects.} A city requires 40\% resident workers on city-funded construction. The city likens its contracts to common property; workers invoke a fundamental common calling. \textit{United Building \& Construction Trades Council v. Camden} (1984) found a fundamental privilege implicated and remanded on justification.
\end{enumerate}

\section{Critique}

\textbf{Natural rights or nondiscrimination?} Washington's paragraph is ambiguous on whether Article IV guarantees substantive natural rights or only interstate equality in fundamental rights. Barnett, Epstein, and Eastman read it as the natural-rights foundation the Reconstruction framers built into the Fourteenth Amendment, making \textit{Slaughter-House} unfaithful; Lash and Berger answer that Article IV was a comity clause, as \textit{Paul v. Virginia} confirmed. Textualists note the anomaly of listing the elective franchise, which visitors have never enjoyed.

\textbf{Vagueness and common property.} Appeals to rights of all free governments supply no clear boundary, which led \textit{Baldwin} to adopt a functional test, and \textit{Toomer} and \textit{Hicklin v. Orbeck} recognized that the state-ownership fiction can mask protectionism. Magliocca notes that the opinion's canonical status was largely a post-Civil War creation.

\section{Key Quotations}

\begin{quote}
``Protection by the government; the enjoyment of life and liberty, with the right to acquire and possess property of every kind, and to pursue and obtain happiness and safety; subject nevertheless to such restraints as the government may justly prescribe for the general good of the whole.''
\end{quote}
\noindent --- Circuit Justice Washington, 6 F. Cas. at 551--52.

\begin{quote}
``The inquiry is, what are the privileges and immunities of citizens in the several states? We feel no hesitation in confining these expressions to those privileges and immunities which are, in their nature, fundamental; which belong, of right, to the citizens of all free governments; and which have, at all times, been enjoyed by the citizens of the several states which compose this Union, from the time of their becoming free, independent, and sovereign.''
\end{quote}
\noindent --- Circuit Justice Washington, 6 F. Cas. at 551.

\begin{quote}
``The right of a citizen of one state to pass through, or to reside in any other state, for purposes of trade, agriculture, professional pursuits, or otherwise; to claim the benefit of the writ of habeas corpus; to institute and maintain actions of any kind in the courts of the state; to take, hold and dispose of property, either real or personal; and an exemption from higher taxes or impositions than are paid by the other citizens of the state; may be mentioned as some of the particular privileges and immunities of citizens, which are clearly embraced by the general description of privileges deemed to be fundamental: to which may be added, the elective franchise, as regulated and established by the laws or constitution of the state in which it is to be exercised.''
\end{quote}
\noindent --- Circuit Justice Washington, 6 F. Cas. at 552.

\end{document}
